% Apunts de teoria en PDF · 1r ESO · Unitat 9: Gràfics i funcions
% Mateix estil que els apunts de Batxillerat (franja i caixes de teoria en blau, exemples en roig), però més visuals,
% curts i amb dibuixos, per a l'alumnat de 3r ESO. Es compila amb pdflatex dues vegades (per als enllaços de l'índex),
% amb logo-fentmates.pdf (_eines/logo/pdf/horitzontal-clar.pdf) al costat.
\documentclass[11pt,a4paper]{article}
\usepackage[utf8]{inputenc}
\usepackage[T1]{fontenc}
\usepackage[catalan]{babel}
\usepackage[sfdefault,book,lining]{FiraSans}
\usepackage[italic,noplusnominus,noequal]{mathastext}
\linespread{1.18}
\usepackage[a4paper,top=2.4cm,bottom=2.2cm,left=1.6cm,right=1.6cm,headheight=1.1cm,headsep=0.5cm,footskip=1cm]{geometry}
\usepackage{amsmath,amssymb,cancel}
\usepackage{icomma}
\usepackage{graphicx}
\usepackage[table]{xcolor}
\usepackage{tikz}\usetikzlibrary{calc,positioning,arrows.meta,decorations.pathreplacing,shapes.geometric}
\usepackage{fontawesome5}
\usepackage{ccicons}
\usepackage{fancyhdr}
\usepackage{enumitem}
\usepackage{array}
\usepackage{needspace}
\usepackage[most]{tcolorbox}
\usepackage{titlesec}
\usepackage[hidelinks]{hyperref}
\hypersetup{pdftitle={Apunts de teoria. Unitat 9: Gràfics i funcions. 1r ESO},pdfsubject={Matemàtiques 1r ESO},pdfkeywords={pla cartesià, eixos de coordenades, quadrants, coordenades d'un punt, plànols, translacions, simetries, magnituds, taules de valors, gràfics, interpretació de gràfics, proporcionalitat directa, funció lineal, pendent}}

% colors de la web (bloc d'Àlgebra: blau)
\definecolor{fmnavy}{HTML}{1B2748}
\definecolor{fmblue}{HTML}{3D6AA0}
\definecolor{fmbluel}{HTML}{4E81BD}
\definecolor{fmyellow}{HTML}{FFC94D}
\definecolor{fmgold}{HTML}{E9A919}
\definecolor{fmaccent}{HTML}{FF7A50}
\definecolor{fmdef}{HTML}{1D4F73}\definecolor{fmdefbg}{HTML}{EEF4F9}
\definecolor{fmex}{HTML}{A3203A}\definecolor{fmexbg}{HTML}{FBF0F2}
\definecolor{fmsf}{HTML}{C2410C}
\definecolor{fmok}{HTML}{15803D}
\definecolor{fmmuted}{HTML}{5D6685}
\definecolor{fmsurface}{HTML}{EEF1F8}
\definecolor{fmA}{HTML}{2F6FD6}\definecolor{fmB}{HTML}{EA6A10}\definecolor{fmR}{HTML}{E11D48}\definecolor{fmK}{HTML}{16A34A}\definecolor{fmV}{HTML}{7C3AED}

% lletres
\newcommand{\comic}{\fontfamily{ComicNeue-TLF}\selectfont}
\newcommand{\heavy}{\fontseries{eb}\selectfont}
\newcommand{\semi}{\fontseries{sb}\selectfont}
\DeclareUnicodeCharacter{2212}{-}

\setlength{\parindent}{0pt}
\setlength{\parskip}{7pt}
\setlength{\lineskip}{4pt}\setlength{\lineskiplimit}{4pt}
\setlength{\emergencystretch}{3em}
\setlist{itemsep=3pt,topsep=3pt,leftmargin=1.4em}
\setcounter{secnumdepth}{1}
\renewcommand{\thesection}{9.\arabic{section}}
\titleformat{\section}[hang]{\color{fmnavy}\fontsize{23}{27}\heavy\selectfont}{\tcbox[on line,colback=fmblue,colframe=fmblue,arc=3mm,boxsep=2pt,left=6pt,right=6pt,top=3pt,bottom=3pt]{\color{white}\thesection}}{0.45em}{}[\vspace{5pt}{\color{fmblue}\titlerule[2pt]}]
\titlespacing*{\section}{0pt}{4pt}{6pt}
\titleformat{\subsection}{\color{fmblue}\Large\semi}{}{0pt}{\faCaretRight\ }
\titlespacing*{\subsection}{0pt}{8pt}{2pt}

% frase de cada apartat, «a mà»
\newcommand{\frase}[1]{\par\vspace{-2pt}{\comic\Large\color{fmbluel}\faCommentDots[regular]\ \ #1\par}\vspace{2pt}}
\newcommand{\seccio}[3]{\Needspace{16\baselineskip}\section{#2}\hypertarget{#1}{}\frase{#3}}
% notes a mà dins dels exemples
\newcommand{\nota}[1]{{\comic\small\color{fmex!85!black}#1}}
\newcommand{\notab}[1]{{\comic\small\color{fmblue}#1}}
% bé / malament
\newcommand{\be}[1]{{\color{fmok}\faCheckCircle}\ \ensuremath{#1}}
\newcommand{\mal}[1]{{\color{fmR}\faTimesCircle}\ {\color{fmR}\ensuremath{#1}}}
% pas numerat
\newcommand{\pas}[1]{\tikz[baseline=(p.base)]\node[circle,fill=fmblue,text=white,inner sep=1.2pt,minimum size=1.25em,font=\small\heavy](p){#1};}
% etiqueta de color
\newcommand{\chip}[2]{\tcbox[on line,colback=#1!14,colframe=#1!60,arc=2mm,boxsep=1pt,left=5pt,right=5pt,top=1pt,bottom=1pt]{#2}}
\newcommand{\fmkey}[1]{\colorbox{fmyellow!55}{\strut #1}}
\newcommand{\cA}[1]{{\color{fmA}#1}}\newcommand{\cB}[1]{{\color{fmB}#1}}\newcommand{\cK}[1]{{\color{fmK}#1}}\newcommand{\cR}[1]{{\color{fmR}#1}}\newcommand{\cV}[1]{{\color{fmV}#1}}

% caixes
\tcbset{fmbase/.style={enhanced,arc=3mm,boxrule=1.2pt,left=4mm,right=4mm,top=2mm,bottom=2.5mm,before skip=8pt,after skip=8pt,fonttitle=\heavy\large,before upper=\raggedright}}
\newtcolorbox{idea}[1][Idea clau]{fmbase,colback=fmdefbg,colframe=fmdef,coltitle=white,title={\faKey\ \ #1}}
\newtcolorbox{exemple}[1][Exemple]{fmbase,colback=fmexbg,colframe=fmex,colbacktitle=fmexbg,coltitle=fmex,title={\faPencil*\ \ #1},titlerule=0.5pt,titlerule style={fmex!40,dashed}}
\newtcolorbox{compte}[1][Compte!]{fmbase,colback=fmsf!6,colframe=fmsf,coltitle=white,fonttitle=\comic\bfseries\Large,title={\faExclamationTriangle\ \ #1}}
\newtcolorbox{truc}[1][Truc]{fmbase,colback=fmyellow!20,colframe=fmgold,colbacktitle=fmyellow,coltitle=fmnavy,fonttitle=\comic\bfseries\Large,title={\faLightbulb[regular]\ \ #1}}

% dibuix de la targeta de la unitat, en blanc
\newcommand{\fmcard}{\begin{tikzpicture}[x=0.062cm,y=-0.062cm,white,line width=3pt,line cap=round]
  \begin{scope}[transparency group,opacity=0.15]
  \node[anchor=base west,inner sep=0pt,font=\fontsize{28}{28}\selectfont] at (14,44) {$P(3,2)$};
  \draw[line width=2pt] (120,70) -- (235,70) (177,120) -- (177,15);
  \fill (210,45) circle (6); \draw[line width=2pt,dashed] (210,45) -- (210,70) (210,45) -- (177,45);
  \end{scope}
\end{tikzpicture}}
\newcommand{\fmbrand}{\raisebox{-0.25\height}{\includegraphics[height=0.5cm]{logo-fentmates.pdf}}}

% regla de Ruffini: #1 = a, #2 = coeficients, #3 = fila del mig (sense el primer), #4 = fila de baix, #5 = nombre de coeficients
\newcommand{\ruffini}[5]{\begin{tikzpicture}[x=1.25cm,y=0.85cm,every node/.style={font=\large}]
  \foreach \v [count=\i] in {#2} \node (t\i) at (\i,0) {$\v$};
  \foreach \v [count=\i] in {#3} {\pgfmathtruncatemacro\j{\i+1}\node (m\j) at (\j,-1) {$\v$};}
  \foreach \v [count=\i] in {#4} \node (b\i) at (\i,-2.1) {$\v$};
  \node[text=fmB] at (0,-1) {$#1$};
  \draw[line width=1.1pt,fmnavy] (0.45,0.5)--(0.45,-1.5) (0.45,-1.5)--(#5+0.5,-1.5);
  \node[draw=fmR,line width=1.3pt,rounded corners=2mm,minimum size=0.75cm] at (#5,-2.1) {};
  \draw[-{Latex[length=2.2mm]},fmK,line width=1.1pt] (t1.south) -- (b1.north);
  \pgfmathtruncatemacro\last{#5-1}
  \foreach \i in {1,...,\last} {\pgfmathtruncatemacro\j{\i+1}
    \draw[-{Latex[length=2.2mm]},fmB,line width=1.1pt,shorten >=1pt] (b\i.north east) -- (m\j.south west);
    \draw[-{Latex[length=2.2mm]},fmA!70,line width=1pt,dashed,shorten >=2pt] ([xshift=8pt]t\j.south) -- ([xshift=8pt]b\j.north);}
\end{tikzpicture}}

% capçalera i peu
\pagestyle{fancy}\fancyhf{}
\fancyhead[L]{\small\fmbrand}
\fancyhead[R]{\small\color{fmmuted}1r ESO · Unitat 9: Gràfics i funcions}
\fancyfoot[L]{\scriptsize\color{fmmuted}\href{https://creativecommons.org/licenses/by-nc-sa/4.0/deed.ca}{\ccbyncsa}\ \ FentMates · fentmates.org · CC BY-NC-SA 4.0}
\fancyfoot[R]{\small\color{fmmuted}\thepage}
\renewcommand{\headrulewidth}{0.4pt}\renewcommand{\headrule}{\color{fmmuted!40}\hrule height \headrulewidth}
\fancypagestyle{primera}{\fancyhf{}\renewcommand{\headrulewidth}{0pt}
  \fancyfoot[L]{\scriptsize\color{fmmuted}\href{https://creativecommons.org/licenses/by-nc-sa/4.0/deed.ca}{\ccbyncsa}\ \ FentMates · fentmates.org · CC BY-NC-SA 4.0}
  \fancyfoot[R]{\small\color{fmmuted}\thepage}}

% targeta de l'índex
\newcommand{\itile}[4]{\begin{tcolorbox}[enhanced,colback=white,colframe=fmblue!35,arc=3mm,boxrule=1pt,boxsep=0pt,left=6pt,right=6pt,top=5pt,bottom=5pt,before skip=0pt,after skip=7pt]
  \hyperlink{#1}{\begin{minipage}[c]{0.9cm}\centering{\color{fmblue}\LARGE #3}\end{minipage}\hspace{6pt}%
  \begin{minipage}[c]{0.78\linewidth}{\color{fmblue}\heavy #2}\\[-2pt]{\color{fmnavy}\semi\large #4}\end{minipage}}\end{tcolorbox}}

% taula amb vores arrodonides: el contingut es retalla amb la forma (sense @{} a les vores ni entre columnes)
\newsavebox\tabclipbox
\newcommand{\taulaclip}[3][3mm]{\sbox\tabclipbox{#3}\begin{tikzpicture}[baseline=(tcn.base)]
  \begin{scope}\clip[rounded corners=#1] (0,-\dp\tabclipbox) rectangle (\wd\tabclipbox,\ht\tabclipbox);
  \node[inner sep=0pt,outer sep=0pt,anchor=base west] (tcn) at (0,0) {\usebox\tabclipbox};\end{scope}
  \draw[#2,rounded corners=#1] (0,-\dp\tabclipbox) rectangle (\wd\tabclipbox,\ht\tabclipbox);
\end{tikzpicture}}
\newcommand{\taularod}[3][fmblue]{\taulaclip{#1,line width=1.2pt}{\renewcommand{\arraystretch}{1.5}\begin{tabular}{#2}#3\end{tabular}}}
% balança
\newcommand{\balanca}[3]{\begin{tikzpicture}[scale=#3,transform shape,every node/.style={anchor=center}]
  \fill[fmnavy] (-0.45,-1.9) rectangle (0.45,-1.75);
  \fill[fmnavy!80] (-0.12,-1.75) rectangle (0.12,0);
  \fill[fmgold] (0,0.05) circle (0.13);
  \draw[fmnavy,line width=2.4pt,line cap=round] (-2,0) -- (2,0);
  \foreach \s in {-1,1} {
    \draw[fmmuted,line width=0.8pt] (\s*2,0) -- (\s*2-0.75,-1.05) (\s*2,0) -- (\s*2+0.75,-1.05);
    \fill[fmyellow!70,draw=fmgold,line width=1pt] (\s*2-0.95,-1.05) arc[start angle=180,end angle=360,x radius=0.95,y radius=0.32] -- cycle;}
  \node[font=\Large,fill=fmA!15,draw=fmA,rounded corners=1.5mm,inner sep=3pt] at (-2,-0.68) {$#1$};
  \node[font=\Large,fill=fmB!15,draw=fmB,rounded corners=1.5mm,inner sep=3pt] at (2,-0.68) {$#2$};
\end{tikzpicture}}
% eixos amb quadrícula (dins d'un tikzpicture): xmin, xmax, ymin, ymax
\newcommand{\eixos}[4]{\draw[fmmuted!22,line width=0.4pt] (#1,#3) grid (#2,#4);
  \draw[-{Latex[length=2mm]},fmmuted,line width=0.8pt] (#1-0.2,0)--(#2+0.4,0) node[right,font=\small]{$x$};
  \draw[-{Latex[length=2mm]},fmmuted,line width=0.8pt] (0,#3-0.2)--(0,#4+0.4) node[above,font=\small]{$y$};}
% taula de traduccions (diccionari) amb les vores arrodonides
\newcommand{\dic}[2]{#1 & {\large$#2$}\\}
\newcommand{\dicratio}{0.62}
\newcommand{\dichead}{\begin{tabular}{@{}>{\raggedright}p{\dimexpr\dicratio\linewidth-12pt}>{\centering\arraybackslash}p{\dimexpr\linewidth-\dicratio\linewidth-12pt}@{}}\faComment[regular]\ Ho diem així\dots & \dots\ així\end{tabular}}
\newenvironment{dicbox}{\begin{tcolorbox}[enhanced,arc=3mm,boxrule=1.2pt,colframe=fmblue,colback=white,coltitle=white,fonttitle=\heavy,
  title={\dichead},left=0pt,right=0pt,top=0pt,bottom=1.5mm,boxsep=0pt,lefttitle=6pt,righttitle=6pt,before skip=0pt,after skip=0pt]\rowcolors{1}{white}{fmsurface}%
  \begin{tabular}{@{\hspace{6pt}}>{\raggedright}p{\dimexpr\dicratio\linewidth-12pt}>{\centering\arraybackslash}p{\dimexpr\linewidth-\dicratio\linewidth-12pt}@{\hspace{6pt}}}}
  {\end{tabular}\end{tcolorbox}}
% resum
\newcommand{\rtile}[3]{\begin{tcolorbox}[enhanced,colback=white,colframe=fmblue!45,arc=3mm,boxrule=1pt,left=3mm,right=3mm,top=2.5mm,bottom=2.5mm,fontupper=\large,
  title={\tcbox[on line,colback=fmblue,colframe=fmblue,arc=2mm,boxsep=0pt,left=4pt,right=4pt,top=1pt,bottom=1pt]{\color{white}\small\heavy #1}\ \ #2},coltitle=fmnavy,colbacktitle=fmdefbg,fonttitle=\semi,toptitle=0.5mm,bottomtitle=0.5mm]#3\end{tcolorbox}}
% portada: #1 subtítol a mà
\newcommand{\portada}[1]{\thispagestyle{primera}
\begin{tikzpicture}[remember picture,overlay]
  \fill[fmblue] (current page.north west) rectangle ([yshift=-9.4cm]current page.north east);
  \begin{scope}
    \clip (current page.north west) rectangle ([yshift=-9.4cm]current page.north east);
    \node[anchor=north west,inner sep=0pt] at ([xshift=6.6cm,yshift=-1.3cm]current page.north west) {\fmcard};
  \end{scope}
  \node[anchor=north west,inner sep=0pt] at ([xshift=1.6cm,yshift=-1.1cm]current page.north west) {\tcbox[colback=white,colframe=white,arc=2.5mm,boxsep=0pt,left=7pt,right=7pt,top=5pt,bottom=5pt]{\includegraphics[height=0.75cm]{logo-fentmates.pdf}}};
\end{tikzpicture}
\vspace*{0.55cm}
\tcbox[colback=white!15!fmblue,colframe=white!15!fmblue,arc=2.5mm,boxsep=0pt,left=6pt,right=6pt,top=3pt,bottom=3pt]{\color{white}\small\bfseries Apunts de teoria · 1r ESO · Matemàtiques}

\vspace{2pt}
{\raggedright\color{white}\fontsize{30}{36}\heavy\selectfont{\color{fmyellow}Unitat 9}\\[2pt]Gràfics i funcions\par}
\vspace{8pt}
{\color{white}\comic\Large #1\par}
\vspace*{\dimexpr 7.4cm-\pagetotal\relax}
}
\newcommand{\llegenda}{\par\vspace{16pt}
{\color{fmnavy}\heavy\Large\faPalette\ \ Com es llegeixen aquests apunts}\par\smallskip
\begin{tcbraster}[raster columns=4,raster equal height,raster column skip=8pt,raster before skip=12pt]
\tcbox[colback=fmdefbg,colframe=fmdef,arc=3mm,boxrule=1.2pt,halign=center,width=\linewidth]{{\color{fmdef}\Large\faKey}\\[2pt]{\color{fmdef}\heavy Idea clau}\\[-1pt]{\small el que cal saber}}
\tcbox[colback=fmexbg,colframe=fmex,arc=3mm,boxrule=1.2pt,halign=center,width=\linewidth]{{\color{fmex}\Large\faPencil*}\\[2pt]{\color{fmex}\heavy Exemple}\\[-1pt]{\small pas a pas}}
\tcbox[colback=fmsf!6,colframe=fmsf,arc=3mm,boxrule=1.2pt,halign=center,width=\linewidth]{{\color{fmsf}\Large\faExclamationTriangle}\\[2pt]{\color{fmsf}\heavy Compte!}\\[-1pt]{\small errors típics}}
\tcbox[colback=fmyellow!20,colframe=fmgold,arc=3mm,boxrule=1.2pt,halign=center,width=\linewidth]{{\color{fmgold!80!black}\Large\faLightbulb[regular]}\\[2pt]{\color{fmnavy}\heavy Truc}\\[-1pt]{\small per a recordar-ho}}
\end{tcbraster}}
\newcommand{\resumtitol}{\clearpage
\hypertarget{resum}{}
{\color{fmnavy}\fontsize{23}{27}\heavy\selectfont\faClipboardList\ \ En resum}\par\vspace{5pt}{\color{fmblue}\titlerule[2pt]}
\frase{Tota la unitat en una pàgina.}}
% els 5 passos per a resoldre problemes: #2 = text del pas 2
\newcommand{\passosproblema}[2]{\begin{center}
\begin{tikzpicture}[s/.style={draw=fmblue,fill=fmdefbg,rounded corners=3mm,line width=1.2pt,minimum width=2.75cm,minimum height=2.45cm,align=center,font=\small},
  ic/.style={font=\LARGE,text=fmblue}]
  \foreach \i/\ic/\t in {1/\faBookOpen/{\textbf{Llegim}\\i entenem},2/\faQuestionCircle/{#1},3/\faEquals/{#2},4/\faCogs/{\textbf{Resolem}},5/\faCheckDouble/{\textbf{Comprovem}\\i responem}} {
    \node[s] (s\i) at (3.3*\i,0) {};
    \node[ic] at ([yshift=-0.5cm]s\i.north) {\ic};
    \node[align=center,font=\small,anchor=north] at ([yshift=-0.95cm]s\i.north) {\t};
    \node[circle,fill=fmblue,text=white,font=\small\heavy,inner sep=1.5pt,minimum size=1.3em] at (s\i.north west) {\i};}
  \foreach \i [evaluate=\i as \j using int(\i+1)] in {1,...,4} \draw[-{Latex[length=3mm]},fmmuted,line width=1.5pt] (s\i) -- (s\j);
\end{tikzpicture}
\end{center}}

\newcommand{\lo}{\mathopen{]}}\newcommand{\ro}{\mathclose{[}}
% punts ple i buit, nombres dels eixos i taula de valors
\newcommand{\pple}[2][fmblue]{\fill[#1] #2 circle (3.6pt);}
\newcommand{\nums}[4]{\foreach \k in {#1,...,#2} {\ifnum\k=0\else\node[below,font=\scriptsize,text=fmmuted] at (\k,0) {$\k$};\fi}
  \foreach \k in {#3,...,#4} {\ifnum\k=0\else\node[left,font=\scriptsize,text=fmmuted] at (0,\k) {$\k$};\fi}}
\newcommand{\tv}[3][fmblue]{\taulaclip{#1,line width=1.2pt}{\renewcommand{\arraystretch}{1.45}%
  \begin{tabular}{>{\columncolor{#1}\color{white}\heavy}c*{12}{w{c}{0.8cm}}}#2\\\rowcolor{#1!8}\cellcolor{#1}#3\end{tabular}}}
% etiqueta d'una recta o corba
\newcommand{\etq}[3]{\node[fill=white,inner sep=1.5pt,text=#1,font=\small] at #2 {#3};}


\newcommand{\passos}[1]{\begin{tcolorbox}[fmbase,colback=fmdefbg,colframe=fmdef,coltitle=white,title={\faListOl\ \ Com es resolen}]#1\end{tcolorbox}}
\def\intvdot#1#2#3{\ifx#2o\filldraw[fill=white,draw=#3,line width=1.3pt] (#1,0) circle (3.3pt);\fi\ifx#2c\fill[#3] (#1,0) circle (3.5pt);\fi}
\newcommand{\intv}[5]{\begin{tikzpicture}[x=0.4cm,baseline=-3pt]
  \draw[fmmuted,line width=0.8pt] (-6,0)--(7,0);
  \foreach \k in {-5,...,6} \draw[fmmuted,line width=0.6pt] (\k,-0.08)--(\k,0.08);
  \ifx#3a\def\xa{-6}\else\def\xa{#1}\fi
  \ifx#4a\def\xb{7}\else\def\xb{#2}\fi
  \ifx#4a\draw[#5,line width=2.6pt,-{Latex[length=2.6mm]}] (\xa,0)--(\xb,0);\else\ifx#3a\draw[#5,line width=2.6pt,{Latex[length=2.6mm]}-] (\xa,0)--(\xb,0);\else\draw[#5,line width=2.6pt] (\xa,0)--(\xb,0);\fi\fi
  \intvdot{#1}{#3}{#5}\intvdot{#2}{#4}{#5}
  \ifx#3a\else\node[below=2pt,font=\scriptsize] at (#1,0) {$#1$};\fi
  \ifx#4a\else\node[below=2pt,font=\scriptsize] at (#2,0) {$#2$};\fi
\end{tikzpicture}}
\newcommand{\pbu}[2][fmblue]{\draw[#1,fill=white,line width=1.2pt] #2 circle (3.4pt);}
\newcommand{\miniax}[4]{\draw[-{Latex[length=1.5mm]},fmmuted,line width=0.6pt] (#1,0)--(#2,0); \draw[-{Latex[length=1.5mm]},fmmuted,line width=0.6pt] (0,#3)--(0,#4);}
\newcommand{\passosT}[2]{\begin{tcolorbox}[fmbase,colback=fmdefbg,colframe=fmdef,coltitle=white,title={\faListOl\ \ #1}]#2\end{tcolorbox}}
\newcommand{\pbuit}[2][fmblue]{\filldraw[fill=white,draw=#1,line width=1.3pt] #2 circle (3.4pt);}
\newcommand{\Dom}{\mathrm{Dom}\,}\newcommand{\Rec}{\mathrm{Rec}\,}


\begin{document}
\portada{Situem punts en el pla i llegim gràfiques.}

{\color{fmnavy}\heavy\Large\faMapSigns\ \ Què aprendràs?}\par\smallskip
\begin{minipage}[t]{0.485\linewidth}
\itile{sec1}{9.1}{\faCrosshairs}{El pla cartesià}
\itile{sec2}{9.2}{\faMapMarker*}{Aplicacions de les coordenades}
\itile{sec3}{9.3}{\faTable}{Magnituds: taules i gràfics}
\itile{sec4}{9.4}{\faBookOpen}{Interpretació de gràfics}
\end{minipage}\hfill
\begin{minipage}[t]{0.485\linewidth}
\itile{sec5}{9.5}{\faChartLine}{Proporcionalitat directa}
\itile{sec6}{9.6}{\faIcon{level-up-alt}}{Funció lineal}
\itile{sec7}{9.7}{\faLightbulb}{Problemes amb funcions lineals}
\itile{resum}{\phantom{9.7}}{\faClipboardList}{En resum}
\end{minipage}
\llegenda

% ============================================================ 9.1
\clearpage
\seccio{sec1}{El pla cartesià. Coordenades d'un punt}{Dos nombres per a dir on és un punt.}

\vfill
\begin{minipage}[c]{0.5\linewidth}
\begin{idea}[El pla cartesià]
Dues rectes numèriques perpendiculars, els \textbf{eixos}:\\[3pt]
\textbf{Eix d'abscisses} ($X$): l'horitzontal; positius a la dreta.\\[2pt]
\textbf{Eix d'ordenades} ($Y$): el vertical; positius amunt.\\[3pt]
Es tallen en l'\textbf{origen} $O$ i fan quatre \textbf{quadrants}: I, II, III i IV, al contrari de les agulles del rellotge.
\end{idea}
\end{minipage}\hfill
\begin{minipage}[c]{0.47\linewidth}\centering
\begin{tikzpicture}[scale=0.5]
  \draw[fmmuted!25] (-5.5,-4.5) grid (5.5,4.5);
  \draw[-{Latex[length=2mm]},fmnavy,line width=1pt] (-5.5,0) -- (5.8,0) node[right,font=\small] {$X$};
  \draw[-{Latex[length=2mm]},fmnavy,line width=1pt] (0,-4.5) -- (0,4.8) node[above,font=\small] {$Y$};
  \foreach \k in {-5,...,-1,1,2,...,5} \node[below,font=\tiny,text=fmmuted] at (\k,0) {$\k$};
  \foreach \k in {-4,...,-1,1,2,...,4} \node[left,font=\tiny,text=fmmuted] at (0,\k) {$\k$};
  \node[font=\large\heavy,text=fmK] at (2.5,2.5) {I}; \node[font=\large\heavy,text=fmA] at (-2.5,2.5) {II};
  \node[font=\large\heavy,text=fmR] at (-2.5,-2.5) {III}; \node[font=\large\heavy,text=fmB] at (2.5,-2.5) {IV};
  \fill[fmnavy] (0,0) circle (3pt) node[below left,font=\small] {$O$};
\end{tikzpicture}
\end{minipage}

\vfill
\begin{minipage}[c]{0.47\linewidth}\centering
\begin{tikzpicture}[scale=0.5]
  \draw[fmmuted!25] (-5.5,-4.5) grid (5.5,4.5);
  \draw[-{Latex[length=2mm]},fmnavy,line width=1pt] (-5.5,0) -- (5.8,0) node[right,font=\small] {$X$};
  \draw[-{Latex[length=2mm]},fmnavy,line width=1pt] (0,-4.5) -- (0,4.8) node[above,font=\small] {$Y$};
  \foreach \k in {-5,...,-1,1,2,...,5} \node[below,font=\tiny,text=fmmuted] at (\k,0) {$\k$};
  \foreach \k in {-4,...,-1,1,2,...,4} \node[left,font=\tiny,text=fmmuted] at (0,\k) {$\k$};
  \draw[fmB,line width=1.4pt,-{Latex[length=2mm]}] (0,0) -- (3,0); \draw[fmK,line width=1.4pt,-{Latex[length=2mm]}] (3,0) -- (3,2);
  \fill[fmR] (3,2) circle (4pt) node[above right,font=\small\heavy] {$A(3,2)$};
  \fill[fmA] (-2,4) circle (4pt) node[right,font=\scriptsize] {$B$}; \fill[fmA] (-4,-1) circle (4pt) node[below,font=\scriptsize] {$C$};
  \fill[fmA] (2,-3) circle (4pt) node[right,font=\scriptsize] {$D$}; \fill[fmV] (4,0) circle (4pt) node[above,font=\scriptsize] {$E$}; \fill[fmV] (0,-2) circle (4pt) node[right,font=\scriptsize] {$F$};
\end{tikzpicture}
\end{minipage}\hfill
\begin{minipage}[c]{0.5\linewidth}
\begin{idea}[Coordenades $P(x,\ y)$]
\textbf{Abscissa} $x$: quant ens movem en \textbf{horitzontal} (dreta $+$, esquerra $-$).\\[2pt]
\textbf{Ordenada} $y$: després, en \textbf{vertical} (amunt $+$, avall $-$).\\[3pt]
$A(3,2)$: 3 a la dreta i 2 amunt.
\end{idea}
\begin{compte}[L'ordre importa]
$(3,2)$ i $(2,3)$ són punts diferents. Primer caminem pel terra ($x$), després pugem per l'escala ($y$).
\end{compte}
\end{minipage}

\vfill
\begin{center}\small
\taularod{>{\heavy}l*{4}{>{\hspace{8pt}}c<{\hspace{8pt}}}}{
\rowcolor{fmblue}\color{white}Quadrant & \color{white}I & \color{white}II & \color{white}III & \color{white}IV\\
Signes & $(+,+)$ & $(-,+)$ & $(-,-)$ & $(+,-)$\\
\rowcolor{fmsurface}Exemple & $A(3,2)$ & $B(-2,4)$ & $C(-4,-1)$ & $D(2,-3)$
}
\end{center}

\vfill
\begin{truc}[Punts sobre els eixos]
Ordenada 0: sobre l'eix $X$, \ $E(4,0)$. \ Abscissa 0: sobre l'eix $Y$, \ $F(0,-2)$. \ L'origen és $O(0,0)$.
\end{truc}

\vfill
% ============================================================ 9.2
\clearpage
\seccio{sec2}{Aplicacions de les coordenades}{Plànols, jocs, figures i moviments.}

\vfill
\begin{minipage}[c]{0.48\linewidth}\centering
\begin{tikzpicture}[scale=0.48]
  \draw[fmmuted!30] (-5.5,-4.5) grid (5.5,4.5);
  \draw[fmnavy,line width=1pt] (-5.5,0) -- (5.5,0) (0,-4.5) -- (0,4.5);
  \node[font=\small] at (0,0) {\faIcon{dot-circle}};
  \node[text=fmA,font=\large] at (1,3) {\faSchool}; \node[text=fmB,font=\large] at (-3,2) {\faHome};
  \node[text=fmK,font=\large] at (4,-2) {\faFutbol}; \node[text=fmK!70!black,font=\large] at (-2,-3) {\faTree};
  \node[text=fmV,font=\large] at (3,1) {\faBook};
\end{tikzpicture}
\end{minipage}\hfill
\begin{minipage}[c]{0.49\linewidth}
\begin{idea}[On és?]
Amb l'origen en la plaça i cada quadre una illa de cases:\\[3pt]
\faSchool\ institut $(1,3)$ \qquad \faHome\ casa $(-3,2)$\\[2pt]
\faFutbol\ poliesportiu $(4,-2)$ \quad \faTree\ parc $(-2,-3)$\\[2pt]
\faBook\ biblioteca $(3,1)$
\end{idea}
\begin{truc}[En la vida real]
Enfonsar la flota i els escacs (B4, E7), la fila i el seient del cine, el GPS: València, $(39^\circ\text{N},\ 0{,}4^\circ\text{O})$.
\end{truc}
\end{minipage}

\vfill
\begin{exemple}[Exemple: \ rectangle $A(-3,-1)$, $B(4,-1)$, $C(4,3)$, $D(-3,3)$]
\begin{minipage}[c]{0.45\linewidth}\centering
\begin{tikzpicture}[scale=0.42]
  \draw[fmmuted!30] (-4,-2) grid (5,4); \draw[fmnavy] (-4,0) -- (5,0) (0,-2) -- (0,4);
  \filldraw[fill=fmA!20,draw=fmA,line width=1.4pt] (-3,-1) rectangle (4,3);
  \foreach \p/\n/\o in {{(-3,-1)}/A/below left,{(4,-1)}/B/below right,{(4,3)}/C/above right,{(-3,3)}/D/above left} \fill[fmA] \p circle (4pt) node[\o,font=\scriptsize] {$\n$};
\end{tikzpicture}
\end{minipage}\hfill
\begin{minipage}[c]{0.52\linewidth}
Mateixa $y$: distància, la diferència de les $x$:\\ $AB=4-(-3)=7$\\[3pt]
Mateixa $x$: diferència de les $y$: \ $BC=3-(-1)=4$\\[3pt]
$P=2\cdot7+2\cdot4=\boxed{22}$ \qquad $A=7\cdot4=\boxed{28}$
\end{minipage}
\end{exemple}

\vfill
\begin{tcbraster}[raster columns=3,raster equal height,raster column skip=8pt]
\begin{tcolorbox}[fmbase,colback=fmK!8,colframe=fmK,coltitle=white,title={Desplaçar}]
3 a la dreta i 2 avall:\par\smallskip
$(x,y)\to(x+3,\ y-2)$\par\smallskip
$(1,4)\to(4,2)$
\end{tcolorbox}
\begin{tcolorbox}[fmbase,colback=fmA!8,colframe=fmA,coltitle=white,title={Simètric respecte de $X$}]
Canvia el signe de la $y$:\par\smallskip
$(x,y)\to(x,\ -y)$\par\smallskip
$(2,3)\to(2,-3)$
\end{tcolorbox}
\begin{tcolorbox}[fmbase,colback=fmB!8,colframe=fmB,coltitle=white,title={Simètric respecte de $Y$}]
Canvia el signe de la $x$:\par\smallskip
$(x,y)\to(-x,\ y)$\par\smallskip
$(2,3)\to(-2,3)$
\end{tcolorbox}
\end{tcbraster}

\vfill
\begin{minipage}[c]{0.45\linewidth}\centering
\begin{tikzpicture}[scale=0.42]
  \draw[fmmuted!30] (-5,-4) grid (5,4); \draw[fmnavy,line width=1pt] (-5,0) -- (5,0) (0,-4) -- (0,4);
  \filldraw[fill=fmK!30,draw=fmK,line width=1.3pt] (1,1) -- (4,1) -- (2,3) -- cycle; \node[font=\scriptsize\heavy] at (2.3,1.7) {$T$};
  \filldraw[fill=fmA!25,draw=fmA,line width=1.3pt] (1,-1) -- (4,-1) -- (2,-3) -- cycle; \node[font=\scriptsize\heavy] at (2.3,-1.7) {$T'$};
  \filldraw[fill=fmB!25,draw=fmB,line width=1.3pt] (-1,1) -- (-4,1) -- (-2,3) -- cycle; \node[font=\scriptsize\heavy] at (-2.3,1.7) {$T''$};
\end{tikzpicture}
\end{minipage}\hfill
\begin{minipage}[c]{0.52\linewidth}
\begin{exemple}[Exemple: \ el triangle $T$: $(1,1)$, $(4,1)$, $(2,3)$]
Simètric respecte de $X$, $T'$: \ $(1,-1)$, $(4,-1)$, $(2,-3)$\\[3pt]
Simètric respecte de $Y$, $T''$: \ $(-1,1)$, $(-4,1)$, $(-2,3)$
\end{exemple}
\end{minipage}

\vfill
% ============================================================ 9.3
\clearpage
\seccio{sec3}{Relacions entre magnituds: taules i gràfics}{Enunciat, taula, gràfica o fórmula: quatre maneres de dir el mateix.}

\vfill
{\large Un taxi cobra \textbf{3 €} per pujar-hi i \textbf{1,20 €} per cada quilòmetre.}

\vfill
\begin{tcbraster}[raster columns=2,raster equal height=rows,raster column skip=10pt,raster row skip=8pt]
\begin{tcolorbox}[fmbase,colback=fmdefbg,colframe=fmdef,coltitle=white,title={\pas{1}\ Enunciat}]
La descripció amb paraules: és el text de dalt.\par\smallskip
\notab{Fàcil d'entendre, però no permet calcular ràpid.}
\end{tcolorbox}
\begin{tcolorbox}[fmbase,colback=fmdefbg,colframe=fmdef,coltitle=white,title={\pas{2}\ Taula de valors},before upper=\centering]
\tv{$x$ (km) & 0 & 1 & 2 & 5 & 10}{$y$ (€) & 3 & 4,2 & 5,4 & 9 & 15}
\end{tcolorbox}
\begin{tcolorbox}[fmbase,colback=fmdefbg,colframe=fmdef,coltitle=white,title={\pas{3}\ Gràfica},before upper=\centering]
\begin{tikzpicture}[x=0.42cm,y=0.13cm]
  \draw[fmmuted!22,line width=0.4pt,xstep=1,ystep=3] (0,0) grid (10,15);
  \draw[-{Latex[length=2mm]},fmmuted,line width=0.8pt] (0,0)--(10.8,0) node[right,font=\small]{km};
  \draw[-{Latex[length=2mm]},fmmuted,line width=0.8pt] (0,0)--(0,16.5) node[above,font=\small]{€};
  \foreach \k in {2,4,6,8,10} \node[below,font=\scriptsize,text=fmmuted] at (\k,0) {\k};
  \foreach \k in {3,6,9,12,15} \node[left,font=\scriptsize,text=fmmuted] at (0,\k) {\k};
  \draw[fmblue,line width=1.6pt] (0,3) -- (10,15);
  \foreach \x/\y in {0/3,1/4.2,2/5.4,5/9,10/15} \fill[fmR] (\x,\y) circle (3.2pt);
\end{tikzpicture}
\end{tcolorbox}
\begin{tcolorbox}[fmbase,colback=fmK!8,colframe=fmK,coltitle=white,title={\pas{4}\ Fórmula (expressió algebraica)},before upper=\centering]
{\Large$f(x)=3+1{,}2\,x$}\par\smallskip
\raggedright La més \textbf{completa}: dona el preu de qualsevol trajecte.\par\smallskip
{\centering$f(7)=3+1{,}2\cdot7=\boxed{11{,}40\text{ €}}$\par}
\end{tcolorbox}
\end{tcbraster}

\vfill
\begin{tcbraster}[raster columns=2,raster equal height,raster column skip=10pt]
\begin{tcolorbox}[fmbase,colback=fmA!8,colframe=fmA,coltitle=white,title={\faArrowRight\ \ Imatge: substituïm}]
La imatge de $a$ és $f(a)$. \ Si $f(x)=x^2-3x+1$:\par\smallskip
{\centering$f(-2)=(-2)^2-3\cdot(-2)+1=4+6+1=\boxed{11}$\par}
\end{tcolorbox}
\begin{tcolorbox}[fmbase,colback=fmB!8,colframe=fmB,coltitle=white,title={\faArrowLeft\ \ Antiimatge: resolem}]
Una antiimatge de $k$ és una $x$ amb $f(x)=k$. \ Si $g(x)=2x+5$, l'antiimatge de 11:\par\smallskip
{\centering$2x+5=11\ \to\ 2x=6\ \to\ \boxed{x=3}$\par}
\end{tcolorbox}
\end{tcbraster}

\vfill
\begin{minipage}[t]{0.42\linewidth}
\begin{compte}[Negatius, entre parèntesis]
\be{(-2)^2=4}\qquad \mal{-2^2=-4}\\[2pt]
{\small En substituir, posa sempre els negatius entre parèntesis.}
\end{compte}
\end{minipage}\hfill
\begin{minipage}[t]{0.55\linewidth}
\renewcommand{\dicratio}{0.62}
\begin{dicbox}
\dic{Perímetre d'un quadrat de costat $x$}{P(x)=4x}
\dic{Àrea d'un quadrat de costat $x$}{A(x)=x^2}
\dic{$x$ entrades a 6 € i 2 € de gestió}{f(x)=6x+2}
\dic{Temps per a fer 120 km a $x$ km/h}{t(x)=\tfrac{120}{x}}
\end{dicbox}
\end{minipage}

\vfill
% ============================================================ 9.4
\clearpage
\seccio{sec4}{Interpretació de gràfics}{Explicar amb paraules, i amb unitats, el que diu la gràfica.}

\vfill
\begin{minipage}[c]{0.55\linewidth}
\passosT{Guió per a llegir un gràfic}{\pas{1} \textbf{Variables}: què hi ha en cada eix i en quines unitats. La independent, en l'horitzontal.\\[3pt]
\pas{2} \textbf{Domini i recorregut}: entre quins valors es mou cada una.\\[3pt]
\pas{3} \textbf{Creixement}: on puja, baixa o és constant. Més inclinat, canvi més ràpid.\\[3pt]
\pas{4} \textbf{Màxims i mínims}: els valors extrems i quan passen.\\[3pt]
\pas{5} \textbf{Valors concrets}: imatges i antiimatges.}
\end{minipage}\hfill
\begin{minipage}[c]{0.42\linewidth}
\begin{compte}[No és un dibuix del camí]
En una gràfica distància a casa–temps, un tram que \textbf{baixa} vol dir que ens \textbf{acostem} a casa, no que baixem una costera; un tram \textbf{horitzontal}, que estem \textbf{parats}.
\end{compte}
\end{minipage}

\vfill
\begin{exemple}[Exemple: \ una excursió a peu (distància a casa, en km, segons el temps, en h)]
\begin{minipage}[c]{0.42\linewidth}\centering
\begin{tikzpicture}[x=0.95cm,y=0.33cm]
  \draw[fmmuted!25] (0,0) grid[xstep=1,ystep=2] (6,10);
  \draw[-{Latex[length=2mm]},fmmuted,line width=0.8pt] (0,0) -- (6.5,0) node[right,font=\scriptsize] {$t$ (h)};
  \draw[-{Latex[length=2mm]},fmmuted,line width=0.8pt] (0,0) -- (0,11) node[above,font=\scriptsize] {$d$ (km)};
  \foreach \k in {1,...,6} \node[below,font=\scriptsize] at (\k,0) {\k};
  \foreach \k in {2,4,...,10} \node[left,font=\scriptsize] at (0,\k) {\k};
  \draw[fmblue,line width=1.8pt] (0,0) -- (2,8) -- (3,8) -- (4,10) -- (6,0);
  \foreach \p in {(0,0),(2,8),(3,8),(4,10),(6,0)} \fill[fmblue] \p circle (2.4pt);
  \draw[fmR,dashed] (0,6) -- (4.8,6); \fill[fmR] (1.5,6) circle (2.4pt); \fill[fmR] (4.8,6) circle (2.4pt);
\end{tikzpicture}
\end{minipage}\hfill
\begin{minipage}[c]{0.55\linewidth}
\pas{1} $t$ en hores, $d$ en km. \ \pas{2} $D=[0,6]$, \ $R=[0,10]$.\\[3pt]
\pas{3} S'allunyen en $\lo0,2\ro\cup\lo3,4\ro$; parats a 8 km entre les 2 h i les 3 h; tornen en $\lo4,6\ro$.\\[3pt]
\pas{4} Màxim $(4,10)$: el punt més llunyà. Mínims $(0,0)$ i $(6,0)$: a casa.\\[3pt]
Velocitats: \ $8:2=4$ km/h; \ $2$ km/h \nota{(tram menys inclinat)}; \ tornen a $10:2=5$ km/h.\\[3pt]
\pas{5} A 6 km de casa: a les $6:4=1{,}5$ h, i a la tornada a les $4+\frac45=4{,}8$ h $=$ \boxed{4\text{ h }48\text{ min}}
\end{minipage}
\end{exemple}

\vfill
\begin{exemple}[Exemple: \ el volum d'aigua (L) d'un dipòsit segons el temps (min)]
\begin{minipage}[c]{0.42\linewidth}\centering
\begin{tikzpicture}[x=0.2cm,y=0.012cm]
  \draw[fmmuted!25] (0,0) grid[xstep=5,ystep=50] (25,200);
  \draw[-{Latex[length=2mm]},fmmuted,line width=0.8pt] (0,0) -- (27,0) node[right,font=\scriptsize] {min};
  \draw[-{Latex[length=2mm]},fmmuted,line width=0.8pt] (0,0) -- (0,225) node[above,font=\scriptsize] {L};
  \foreach \k in {5,10,...,25} \node[below,font=\scriptsize] at (\k,0) {\k};
  \foreach \k in {100,200} \node[left,font=\scriptsize] at (0,\k) {\k};
  \draw[fmK!80!black,line width=1.8pt] (0,0) -- (10,200) -- (15,200) -- (25,0);
\end{tikzpicture}
\end{minipage}\hfill
\begin{minipage}[c]{0.55\linewidth}
En $\lo0,10\ro$ s'omple a $200:10=\boxed{20\text{ L/min}}$.\\[3pt]
En $\lo10,15\ro$ es manté amb 200 L (constant).\\[3pt]
En $\lo15,25\ro$ es buida, també a 20 L/min.\\[3pt]
Hi ha 100 L als \boxed{5\text{ min}} (omplint-se) i als \boxed{20\text{ min}} (buidant-se).
\end{minipage}
\end{exemple}

\vfill
\begin{truc}[Com més inclinat, més ràpid]
La rapidesa del canvi és el que puja (o baixa) la $y$ per cada unitat de $x$: en una gràfica distància–temps, és la \textbf{velocitat}.
\end{truc}

\vfill

\vfill
% ============================================================ 9.5
\clearpage
\seccio{sec5}{Funció de proporcionalitat directa}{El doble de $x$, el doble de $y$: una recta per l'origen.}

\vfill
\begin{minipage}[c]{0.5\linewidth}
\begin{idea}[Funció de proporcionalitat directa]
{\centering\fontsize{22}{26}\selectfont$y=\cA{m}\,x$\par}\smallskip
\chip{fmA}{$m$} és la \textbf{constant de proporcionalitat}: \ $m=\dfrac{y}{x}$, igual per a tots els parells.\\[3pt]
La gràfica és una \textbf{recta que passa per l'origen} $(0,\,0)$.\\[3pt]
$m$ també és el \textbf{pendent}: el que augmenta la $y$ quan la $x$ augmenta 1.
\end{idea}
\end{minipage}\hfill
\begin{minipage}[c]{0.47\linewidth}\centering
\begin{tikzpicture}[scale=0.68]
  \eixos{-3}{3}{-3}{3}
  \draw[fmK,line width=1.6pt] (-1.5,-3) -- (1.5,3); \draw[fmK!60,line width=1.6pt] (-3,-1.5) -- (3,1.5);
  \draw[fmR,line width=1.6pt] (-1,3) -- (1,-3); \draw[fmR!60,line width=1.6pt] (-3,1) -- (3,-1);
  \node[above,font=\small,text=fmK!80!black] at (1.5,3) {$y=2x$}; \node[right,font=\small,text=fmK!60!black] at (3,1.5) {$y=\frac12x$};
  \node[above,font=\small,text=fmR] at (-1,3) {$y=-3x$}; \node[left,font=\small,text=fmR!70!black] at (-3,1) {$y=-\frac13x$};
\end{tikzpicture}
\end{minipage}

\vfill
\begin{tcbraster}[raster columns=3,raster equal height,raster column skip=8pt]
\begin{tcolorbox}[fmbase,colback=fmK!8,colframe=fmK,coltitle=white,title={$m>0$},before upper=\centering]
\textbf{Creixent}: la recta puja.
\end{tcolorbox}
\begin{tcolorbox}[fmbase,colback=fmR!6,colframe=fmR,coltitle=white,title={$m<0$},before upper=\centering]
\textbf{Decreixent}: la recta baixa.
\end{tcolorbox}
\begin{tcolorbox}[fmbase,colback=fmdefbg,colframe=fmdef,coltitle=white,title={$|m|$ més gran},before upper=\centering]
Més \textbf{inclinada}.
\end{tcolorbox}
\end{tcbraster}

\vfill
\begin{minipage}[c]{0.58\linewidth}
\begin{truc}[Basta un punt]
L'origen ja el tenim. Per a $y=\tfrac32x$, triem una $x$ que faça la $y$ entera: \ $x=2\ \to\ y=3$. \ Unim $(0,\,0)$ i $(2,\,3)$.
\end{truc}
\end{minipage}\hfill
\begin{minipage}[c]{0.38\linewidth}\centering
\begin{tikzpicture}[scale=0.45]
  \eixos{-2}{4}{-2}{4}
  \draw[fmblue,line width=1.6pt] (-1.33,-2) -- (2.67,4);
  \pple[fmR]{(0,0)} \pple[fmR]{(2,3)}
  \node[right,font=\small\heavy,text=fmR] at (2.1,2.8) {$(2,\,3)$};
\end{tikzpicture}
\end{minipage}

\vfill
\begin{exemple}[Exemple: \ un quilo de cireres costa 4,50 €]
\begin{minipage}[c]{0.5\linewidth}\centering
\tv[fmex]{$x$ (kg) & 0 & 1 & 2 & 3 & 4}{$y$ (€) & 0 & 4,5 & 9 & 13,5 & 18}
\end{minipage}\hfill
\begin{minipage}[c]{0.47\linewidth}
$y=4{,}5\,x$ \ \nota{$m$: el preu d'1 kg}\\[3pt]
$2{,}4$ kg: \ $4{,}5\cdot2{,}4=\boxed{10{,}80\text{ €}}$
\end{minipage}
\end{exemple}

\vfill
\begin{exemple}[Exemple: \ un cotxe gasta 6 L cada 100 km. Quants km fa amb 45 L?]
$m=\dfrac{6}{100}=0{,}06$ L/km \ $\to$ \ $y=0{,}06\,x$ \qquad $0{,}06\,x=45\ \to\ x=\dfrac{45}{0{,}06}=\boxed{750\text{ km}}$
\end{exemple}

\vfill
\begin{exemple}[Exemple: \ passa per $(-4,\,6)$. Passa també per $(10,\,-15)$?]
$m=\dfrac{6}{-4}=-\dfrac32\ \to\ \boxed{y=-\tfrac32x}$ \qquad $x=10:\ \ y=-\tfrac32\cdot10=-15$ \ \be{} \ \ \textbf{sí}
\end{exemple}

\vfill
% ============================================================ 9.6
\clearpage
\seccio{sec6}{Funció lineal}{La recta de $y=mx$, pujada o baixada $n$ unitats.}

\vfill
\begin{center}
\begin{tikzpicture}[t/.style={font=\fontsize{30}{30}\selectfont,inner sep=3pt,text height=25pt,text depth=8pt}]
  \node[t] (y) {$y=$};
  \node[t,right=2pt of y,fill=fmA!12,rounded corners=2mm] (m) {$\cA{m}$};
  \node[t,right=0pt of m] (x) {$x$};
  \node[t,right=2pt of x] (p) {$+$};
  \node[t,right=2pt of p,fill=fmB!12,rounded corners=2mm] (n) {$\cB{n}$};
  \node[fill=fmA!14,draw=fmA,rounded corners=2mm,font=\large,anchor=east,align=center] (lm) at ([xshift=-1.4cm,yshift=0.5cm]y.west) {{\heavy\color{fmA}pendent}\\[-2pt]\small el que varia la $y$ si la $x$ augmenta 1};
  \node[fill=fmB!14,draw=fmB,rounded corners=2mm,font=\large,anchor=west,align=center] (ln) at ([xshift=1.2cm,yshift=0.5cm]n.east) {{\heavy\color{fmB}ordenada en l'origen}\\[-2pt]\small talla l'eix $Y$ en $(0,\,n)$};
  \draw[-{Latex[length=3mm]},fmA,line width=1.3pt] (lm.east) to[out=0,in=110] ([yshift=2pt]m.north);
  \draw[-{Latex[length=3mm]},fmB,line width=1.3pt] (ln.west) to[out=180,in=70] ([yshift=2pt]n.north);
\end{tikzpicture}
\end{center}

\vfill
\begin{minipage}[c]{0.36\linewidth}\centering
\begin{tikzpicture}[scale=0.5]
  \eixos{-3}{3}{-3}{4}
  \draw[fmmuted,line width=1.4pt,dashed] (-1.5,-3) -- (1.5,3);
  \draw[fmblue,line width=1.7pt] (-3,-3) -- (0.5,4);
  \pple[fmB]{(0,3)}
  \draw[-{Latex[length=2mm]},fmB,line width=1.1pt] (0.5,0.5) -- (0.5,2.6);
  \etq{fmblue}{(-2.3,1)}{$y=2x+3$} \etq{fmmuted}{(2.2,1.6)}{$y=2x$}
\end{tikzpicture}\\
{\comic\small\color{fmmuted}la mateixa recta, 3 unitats amunt}
\end{minipage}\hfill
\begin{minipage}[c]{0.6\linewidth}
\begin{idea}[Com es dibuixa]
\pas{1} Marquem $(0,\,n)$.\\[3pt]
\pas{2} Avancem \textbf{1} a la dreta i \textbf{pugem $m$} (o baixem si $m<0$). Si $m$ és una fracció $\tfrac{a}{b}$: avancem $b$ i pugem $a$.\\[3pt]
\pas{3} Unim els dos punts.
\end{idea}
\end{minipage}

\vfill
\begin{minipage}[c]{0.5\linewidth}
\begin{exemple}[Exemple: \ $y=-\tfrac23x+2$]
\begin{minipage}[c]{0.5\linewidth}\centering
\begin{tikzpicture}[scale=0.42]
  \eixos{-2}{5}{-2}{4}
  \draw[fmblue,line width=1.6pt] (-1.5,3) -- (4.5,-1);
  \draw[fmK,line width=1.1pt,-{Latex[length=1.8mm]}] (0,2) -- (3,2);
  \draw[fmR,line width=1.1pt,-{Latex[length=1.8mm]}] (3,2) -- (3,0);
  \pple[fmB]{(0,2)} \pple[fmB]{(3,0)}
\end{tikzpicture}
\end{minipage}\hfill
\begin{minipage}[c]{0.48\linewidth}\small
$n=2$: \ $(0,\,2)$\\[2pt]
\cK{avancem 3}, \cR{baixem 2}:\\ $(3,\,0)$
\end{minipage}
\end{exemple}
\end{minipage}\hfill
\begin{minipage}[c]{0.47\linewidth}
\begin{tcolorbox}[fmbase,colback=fmdefbg,colframe=fmdef,coltitle=white,title={Talls amb els eixos}]
Eix $Y$: \ $x=0\ \to\ (0,\,n)$\\[2pt]
Eix $X$: \ $y=0$, resolem $mx+n=0$\\[4pt]
{\centering$y=3x-6$: \ $(0,\,-6)$ \ i \ $3x-6=0\to(2,\,0)$\par}
\end{tcolorbox}
\end{minipage}

\vfill
\begin{tcbraster}[raster columns=2,raster equal height,raster column skip=10pt]
\begin{tcolorbox}[fmbase,colback=fmsurface,colframe=fmmuted,coltitle=white,title={Funció constant: \ $m=0$}]
\begin{minipage}[c]{0.55\linewidth}$y=n$ és una recta \textbf{horitzontal}.\end{minipage}\hfill
\begin{minipage}[c]{0.4\linewidth}\centering\begin{tikzpicture}[scale=0.3]\eixos{-3}{3}{-2}{3}\draw[fmblue,line width=1.5pt] (-3,2) -- (3,2);\etq{fmblue}{(1.6,2.9)}{$y=2$}\end{tikzpicture}\end{minipage}
\end{tcolorbox}
\begin{tcolorbox}[fmbase,colback=fmR!6,colframe=fmR,coltitle=white,title={Recta vertical: \ $x=k$}]
\begin{minipage}[c]{0.55\linewidth}\textbf{No és una funció}: a $k$ li corresponen infinites $y$.\end{minipage}\hfill
\begin{minipage}[c]{0.4\linewidth}\centering\begin{tikzpicture}[scale=0.3]\eixos{-3}{3}{-2}{3}\draw[fmR,line width=1.5pt] (-2,-2) -- (-2,3);\etq{fmR}{(-1.6,-2.8)}{$x=-2$}\end{tikzpicture}\end{minipage}
\end{tcolorbox}
\end{tcbraster}

\vfill
\begin{compte}[Primer, en la forma $y=mx+n$]
$y=\dfrac{2-3x}{4}=\dfrac24-\dfrac34x\ \to\ y=-\dfrac34x+\dfrac12$: \ $m=-\tfrac34$, \ $n=\tfrac12$. \qquad \mal{m=2}
\end{compte}

\vfill
% ============================================================ 9.7
\clearpage
\seccio{sec7}{Problemes amb funcions lineals}{Una part fixa i una part que creix al mateix ritme.}

\vfill
\begin{idea}
{\centering\Large$y=\cA{m}\,x+\cB{n}$\qquad \normalsize\chip{fmB}{$n$: el valor inicial (quan $x=0$)} \ \ \chip{fmA}{$m$: el ritme de canvi per unitat}\par}
\end{idea}

\vfill
\begin{center}
\taularod{l<{\hspace{6pt}}>{\hspace{6pt}}c<{\hspace{6pt}}>{\hspace{6pt}}c<{\hspace{6pt}}>{\hspace{6pt}}c<{\hspace{4pt}}}{
\rowcolor{fmblue}\color{white}\heavy Situació & \color{white}\heavy Funció & \color{white}\heavy $m$ & \color{white}\heavy $n$\\
Taxi: 3 € fixos i 1,20 € per km & $P=1{,}2x+3$ & preu per km & preu fix\\
\rowcolor{fmsurface}Moviment a velocitat constant & $d=d_0+v\,t$ & velocitat & posició inicial\\
Graus Celsius $\to$ Fahrenheit & $F=1{,}8\,C+32$ & $1{,}8$ & $32$\\
\rowcolor{fmsurface}Dipòsit de 500 L que es buida a 25 L/min & $V=500-25t$ & $-25$ & $500$
}
\end{center}

\vfill
\begin{minipage}[c]{0.45\linewidth}\raggedright
\pas{1} \textbf{Variables}: què és la $x$ i què la $y$ (amb unitats).\\[5pt]
\pas{2} Troba $m$ i $n$. Si et donen \textbf{dos parells}, calcula la recta que passa per aquests punts.\\[5pt]
\pas{3} Escriu la fórmula, respon i \textbf{interpreta}.
\end{minipage}\hfill
\begin{minipage}[c]{0.52\linewidth}
\begin{exemple}[Exemple: \ bicis a 5 € fixos i 2,50 € per hora]
$P(x)=2{,}5x+5$\\[3pt]
4 h: \ $P(4)=10+5=\boxed{15\text{ €}}$\\[3pt]
Amb 25 €: \ $2{,}5x+5=25\ \to\ 2{,}5x=20\ \to\ \boxed{x=8\text{ h}}$
\end{exemple}
\end{minipage}

\vfill
\begin{exemple}[Exemple: \ un rellotge de sorra té 300 g a dalt i en cauen 12 g per minut]
\begin{minipage}[c]{0.55\linewidth}
$S(t)=300-12t$\\[4pt]
Es buida quan $S=0$: \ $12t=300\ \to\ t=25$ min\\[4pt]
Domini amb sentit: \ $\boxed{[0,\,25]}$
\end{minipage}\hfill
\begin{minipage}[c]{0.42\linewidth}\centering
\begin{tikzpicture}[x=0.17cm,y=0.0082cm]
  \draw[fmmuted!22,line width=0.4pt,xstep=5,ystep=50] (0,0) grid (25,300);
  \draw[-{Latex[length=2mm]},fmmuted,line width=0.8pt] (0,0)--(27.5,0) node[right,font=\small]{$t$};
  \draw[-{Latex[length=2mm]},fmmuted,line width=0.8pt] (0,0)--(0,330) node[above,font=\small]{$S$};
  \draw[fmex,line width=1.7pt] (0,300) -- (25,0);
  \pple[fmex]{(0,300)} \pple[fmex]{(25,0)}
  \node[left,font=\scriptsize] at (0,300) {300}; \node[below,font=\scriptsize] at (25,0) {25};
\end{tikzpicture}
\end{minipage}
\end{exemple}

\vfill
\begin{minipage}[t]{0.5\linewidth}
\begin{exemple}[Exemple: \ de °C a °F]
L'aigua bull a 100 °C $=$ 212 °F i es congela a 0 °C $=$ 32 °F:\\[3pt]
$m=\dfrac{212-32}{100-0}=1{,}8$, \ $n=32$ \ $\to$ \ $F=1{,}8\,C+32$
\end{exemple}
\end{minipage}\hfill
\begin{minipage}[t]{0.47\linewidth}
\begin{compte}[Només per a alguns valors]
La fórmula té sentit en un \textbf{domini}: el temps no és negatiu, un dipòsit no pot tindre menys de 0 L\dots
\end{compte}
\end{minipage}

\vfill
% ============================================================ resum
\resumtitol
\begin{tcbraster}[raster columns=2,raster equal height=rows,raster column skip=14pt,raster row skip=34pt,raster before skip=18pt]
\rtile{9.1}{Pla cartesià}{$P(x,y)$: primer l'horitzontal.\\[1ex] Quatre quadrants.}
\rtile{9.2}{Aplicacions}{Plànols i figures.\\[1ex] Simètric respecte de $X$: $(x,-y)$.\\[1ex] Simètric respecte de $Y$: $(-x,y)$.}
\rtile{9.3}{Magnituds}{Independent ($x$) i dependent ($y$).\\[1ex] Taula, gràfic i fórmula.}
\rtile{9.4}{Interpretar}{Variables, domini, creixement,\\[1ex] extrems i valors, amb unitats.}
\rtile{9.5}{Proporcionalitat}{$y=mx$: recta per l'origen.}
\rtile{9.6}{Funció lineal}{$y=mx+n$.\\[1ex] $m$: pendent; $n$: tall amb l'eix $Y$.}
\rtile{9.7}{Problemes}{$n$: valor inicial.\\[1ex] $m$: el que augmenta per unitat.}
\end{tcbraster}
\end{document}
